% Part: methods
% Chapter: proofs
% Section: reading-proofs

\documentclass[../../../include/open-logic-section]{subfiles}

\begin{document}

\olfileid[tr]{mth}{prf}{rea}
\olsection{Kanıtları Okumak}

Ders kitaplarında ve makalelerde bulduğunuz kanıtlar, şimdiye dek
örneklerimize kattığımız bütün ayrıntıları çok ender verir. Yazarlar
çoğu kez durumları ayırdıklarına ya da dolaylı kanıt verdiklerine dikkat
çekmez, hatta bir tanım kullandıklarını belirtmezler. Bu nedenle bir
ders kitabındaki kanıtı okurken onu anlayabilmek için bu ayrıntıları
çoğu zaman kendiniz tamamlamalısınız. Bunu yapmak, bir kanıtta
gerçekleştirmeniz gereken çeşitli hamlelere alışmak için de iyi bir
uygulamadır. Bir örneğe bakalım.

\begin{prop}[Soğurma]
Bütün $A$, $B$ kümeleri için
\[
A \cap (A \cup B) = A
\]
\end{prop}

\begin{proof}
$z \in A \cap (A \cup B)$ ise $z \in A$; dolayısıyla $A \cap (A \cup
B) \subseteq A$. Şimdi $z \in A$ olduğunu varsayalım. O zaman $z \in
A \cup B$ de geçerlidir ve dolayısıyla $z \in A \cap (A \cup B)$ de
geçerlidir.
\end{proof}

Yukarıdaki soğurma yasası kanıtı çok yoğundur. Kullanılan tanımlardan
söz edilmez, bir şeyi kanıtlamadan önce ``şunu kanıtlamalıyız'' denmez
vb. Kanıtı açalım. Kanıtlanan önerme, herhangi $A$ ve $B$ kümeleri
hakkında genel bir iddiadır; kanıtta $A$ ya da $B$ geçtiğinde bunlar
keyfî kümelerin değişkenleridir. Kanıtın ortaya koyduğu genel iddialar,
kümelerin özdeşliğini kanıtlamak için gereken şeydir: yani özdeşliğin
sol tarafındaki her !!{element}, sağ tarafta bir !!a{element} olarak bulunur ve
bunun tersi de geçerlidir.

\begin{quote}
``$z \in A \cap (A \cup B)$ ise $z \in A$; dolayısıyla $A \cap (A \cup
B) \subseteq A$.''
\end{quote}

Bu, özdeşlik kanıtının ilk yarısıdır: keyfî bir~$z$ sol tarafta bir
!!a{element} olarak bulunuyorsa sağ tarafta da bir !!a{element} olarak bulunduğunu, yani $A
  \cap (A \cup B) \subseteq A$ olduğunu belirler. $z \in A \cap (A \cup
  B)$ olduğunu varsayın. $z$'nin iki kümenin kesişiminde bir !!{element}
  olması, ancak ve ancak her iki kümede de bir !!{element} olmasına denk
  olduğundan, $z \in A$ ve aynı zamanda $z \in A \cup B$ sonucuna varabiliriz. Özellikle $z \in
A$'dır; göstermek istediğimiz de buydu. İlk yarı için yapılması gereken
yalnızca bu olduğundan, kanıtın geri kalanının ikinci yarının kanıtı,
yani $A \subseteq A \cap (A \cup B)$'nin bir kanıtı olması gerektiğini
biliyoruz.

\begin{quote}
``Şimdi $z \in A$ olduğunu varsayalım. O zaman $z \in A \cup B$ de
geçerlidir ve dolayısıyla $z \in A \cap (A \cup B)$ de geçerlidir.''
\end{quote}

$z \in A$ olduğunu varsayarak başlarız; çünkü her~$z$ için $z \in A$
ise $z \in A \cap (A \cup B)$ olduğunu gösteriyoruz. $z \in A
\cap (A \cup B)$ olduğunu göstermek için, ``$\cap$'' tanımı gereği (i)
$z \in A$ ve ayrıca (ii) $z \in A \cup B$ olduğunu göstermeliyiz.
Burada (i) zaten varsayımımızdır; dolayısıyla kanıtlanacak başka bir
şey yoktur ve kanıtın bunu yeniden anmamasının nedeni budur. (ii) için,
$z$'nin bir küme birleşiminde !!a{element} olarak bulunmasının ancak ve ancak bu
kümelerden en az birinde !!{element} olarak bulunmasıyla geçerli olduğunu
hatırlayın. $z \in A$ ve $A \cup B$, $A$ ile $B$'nin birleşimi olduğuna
göre burada durum budur. Öyleyse $z \in A \cup B$. Hem (i) $z \in A$
hem de (ii) $z \in A \cup B$ olduğunu gösterdik; dolayısıyla ``$\cap$''
tanımı gereği $z \in A \cap (A \cup B)$. Kanıt bu tanımlardan söz etmez;
okurun onları zaten içselleştirdiği varsayılır. İçselleştirmediyseniz,
geri dönüp ne olduklarını kendinize hatırlatmalısınız. Ardından $z \in
A$'dan neden $z \in A \cup B$'nin ve $z \in A$ ile $z \in A \cup B$'den
neden $z \in A \cap (A \cup B)$'nin çıktığını da görmelisiniz.

Yukarıdaki kanıtın, her şeyin açıkça yazıldığı başka bir biçimi şöyledir:
\begin{proof}{}
[Kümeler için $=$ tanımı gereği $A \cap (A \cup B) = A$'yı kanıtlamak
  üzere (a) $A \cap (A \cup B) \subseteq A$ ve (b) $A \subseteq A \cap
  (A \cup B)$ olduğunu göstermeliyiz. (a): $\subseteq$ tanımı gereği $z
  \in A \cap (A \cup B)$ ise $z \in A$ olduğunu göstermeliyiz.] $z \in
A \cap (A \cup B)$ ise $z \in A$ [$\cap$ tanımı gereği $z \in A \cap
  (A \cup B)$, ancak ve ancak $z \in A$ ve $z \in A \cup B$ ise
  geçerli olduğundan]; dolayısıyla $A \cap (A \cup B) \subseteq A$.
[(b): $\subseteq$ tanımı gereği $z \in A$ ise $z \in A \cap (A \cup
  B)$ olduğunu göstermeliyiz.] Şimdi [(1)] $z \in A$ olduğunu
varsayalım. O zaman [(2)] $z \in A \cup B$ de geçerlidir [(1)'e göre
  $z \in A$ ya da $z \in B$ olduğundan; bu da $\cup$ tanımı gereği $z
  \in A \cup B$ demektir] ve dolayısıyla $z \in A \cap (A \cup B)$ de
geçerlidir [$\cap$ tanımı $z \in A$'yı, yani (1)'i ve $z \in A \cup B$'yi,
  yani (2)'yi gerektirdiğinden].
\end{proof}

\begin{prob}
$A \cup (A \cap B) = A$'nın aşağıdaki kanıtını genişletin: kullanılan
bütün çıkarım kalıplarını, her adımın varsayımlardan ya da daha önce
belirlenmiş iddialardan neden çıktığını ve hangi tanımlara nerede
başvurmak gerektiğini belirtin.
\begin{proof}
  $z \in A \cup (A \cap B)$ ise $z \in A$ ya da $z \in A \cap B$. $z
  \in A \cap B$ ise $z \in A$. Her $z \in A$ aynı zamanda $\in A \cup
  (A \cap B)$'dir.
\end{proof}
\end{prob}

\end{document}
